\documentclass[11pt]{article}
\usepackage[margin=1in]{geometry}
\usepackage{amsmath, amssymb}
\usepackage{booktabs}
\usepackage{graphicx}
\usepackage{hyperref}
\usepackage{xcolor}

\title{Self-Prediction as Information-Bottleneck Regularization in Growing Neural Networks}
\author{Zaccarie Aglas}
\date{May 2026}

\begin{document}
\maketitle

\begin{abstract}
We study adaptive neural networks that grow their own architecture during training (VerifyGrowNet, VG) coupled with a trajectory analysis pipeline (Dynamoscope) for measurement of latent dynamics. Across nine sprints we find: (1) growing network state evolves as a low-rank dynamical system with a universal spectral slope of $-1.85 \pm 0.13$ across seeds; (2) structural signals from dynamoscope ($\varphi_{id}$ emergence, local Lyapunov) outperform loss-based plateau detection as growth triggers on simple tasks; (3) random-init tiny CNNs (498 params) extract trajectories closer to ground truth than ResNet50 (23.5M params) on visual dynamical systems; (4) explicit self-prediction loss as auxiliary objective acts as information-bottleneck regularization, with optimal weight depending on use case: $\alpha \approx 1$ for direct task accuracy, $\alpha \approx 2$ for downstream transfer; (5) hierarchical observer stacks do not produce meta-information --- recursive observation produces redundancy with degraded predictability. These findings constitute a partial vindication and quantification of active inference (Friston) in computational substrates while challenging scale-first inductive bias dogma.
\end{abstract}

\section{Introduction}

Adaptive neural networks --- networks that modify their own architecture during training --- bridge two long-standing intuitions: neuroplasticity in biology and self-organizing systems in physics. The VerifyGrowNet (VG) framework~\cite{aglas2026verifygrow} introduces \emph{discipline} into growth: a mutation is kept only if it produces measurable validation gain. This produces small efficient networks that match or beat fixed architectures.

Dynamoscope~\cite{aglas2026dynamoscope} treats latent trajectories of pretrained encoders applied to video as state-space objects, applying nonlinear dynamics tools (Lyapunov exponents, Koopman modes, persistent homology, SINDy equation discovery, $\varphi_{id}$ causal emergence).

In this work we fuse the two: we treat the growing architecture itself as a dynamical system observable by dynamoscope, then close the loop by using dynamoscope signals as control signals back into VG growth decisions, then close it further with self-observation. We document seven specific findings via empirical sprints on MNIST, CIFAR-100, FashionMNIST, and synthetic Lorenz trajectories.

\textbf{Question structuring the work}: \emph{When a network modifies itself, what laws govern its evolution, and can those laws be used as control signals for the modification?}

\section{Related Work}

\paragraph{Neural Architecture Search (NAS).} ENAS, DARTS, ProxylessNAS optimize architecture via outer loop. VG differs: growth happens \emph{during} a single training, gated by validation gain.

\paragraph{Predictive coding \& active inference.} Rao-Ballard, Friston: minimize prediction error of own model of world. Sprint 5C/5C.1 implements this with the system observing itself.

\paragraph{Information Bottleneck.} Tishby: representation compression maximizes task-irrelevant invariance. We find self-prediction loss acts as similar pressure (Sprint 7, 8B).

\paragraph{Self-organized criticality.} Bak-Tang-Wiesenfeld. Sprint 4 spectral slope $-1.85$ across seeds suggests SOC signature.

\paragraph{Strange loops.} Hofstadter \emph{I Am a Strange Loop}: self-reference produces emergent identity. Sprint 5C/8C test concrete instantiations.

\section{Method}

\subsection{Building blocks}
\begin{itemize}
    \item \textbf{VG-MLP / VG-CNN}: networks with verify-then-keep mutation (widen via GradMax-orthogonal init, function-preserving).
    \item \textbf{Architecture state vector}: 58-dim representation of VG state (layer widths, effective rank, gradient norms, etc.) ; invariant to growth.
    \item \textbf{Architecture image}: render of weight/gradient/bias magnitudes as fixed-size RGB ($64 \times 64$, also invariant).
    \item \textbf{Dynamoscope analyses}: Lyapunov, $\varphi_{id}$, SINDy R$^2$, spectral slope, RQA --- all callable as HTTP service.
\end{itemize}

\subsection{Sprint structure}
We describe seven coupling experiments (Table~\ref{tab:sprints}).

\begin{table}[h]
\caption{Sprint summary. ``Robust'' = replicated on second dataset.}
\label{tab:sprints}
\centering\small
\begin{tabular}{rlll}
\toprule
\# & Title & Finding & Robust ? \\
\midrule
4 & Meta-loop external & spectral\_slope $-1.85$ universal & not retested \\
2 & Reverse trigger & emergence trigger +2.6\% acc on MNIST & $\times$ (CIFAR-100 $\Delta=0.06\%$) \\
3 & Adaptive encoder & vg\_tiny (498p) $6\times$ closer to GT than ResNet50 & directional \\
5B & Visual coupling & visual observer matches direct latent observation & internal \\
5C & Strange loop $\alpha{=}0$ & task converges, self-model diverges & repeatable \\
5C.1 & Self-coh weight sweep & $\alpha{\approx}1$ optimal: cheap regularization & $\checkmark$ (CIFAR-100 $+1.7\%$) \\
7 & Transfer & $\alpha{=}2$ best transfer FashionMNIST $+0.9\%$ & directional, n=3 \\
8B & Multi-step horizon & no effect (state stabilizes too fast) & null finding \\
8C & Hierarchical obs & higher orders diverge, encode redundantly & paradoxical \\
\bottomrule
\end{tabular}
\end{table}

\section{Results}

\subsection{Architecture-as-trajectory (Sprint 4)}
Tracing VG architecture state during MNIST training and feeding into dynamoscope produces measurable dynamics with strikingly consistent inter-seed properties:

\begin{itemize}
    \item Spectral slope: $-1.85 \pm 0.13$ (5 seeds, std 7\%)
    \item DMD rank: 2--3 (system is intrinsically low-rank in 58-D ambient)
    \item RQA DET=LAM=$0.999$ (deterministic + laminar)
    \item $\varphi_{id} = 0.19 \pm 0.29$ (causal emergence detected, "macro more predictable than micro")
    \item Bimodal basin in shared UMAP: intra-cluster distance $\approx 0.05$, inter-cluster $\approx 0.30$ (6$\times$ separation)
\end{itemize}

The $1/f^2$ spectral universality across seeds suggests self-organized criticality.

\subsection{Reverse trigger (Sprint 2)}
Comparing growth triggers on MNIST:
\begin{itemize}
    \item LossPlateau: 0/3 seeds fired, acc $0.876 \pm 0.004$
    \item EmergenceTrigger ($\varphi_{id}$ drop): 3.3 fires per run, acc $\mathbf{0.902 \pm 0.002}$ ($+2.6\%$)
    \item LyapunovLocalTrigger: 3.3 fires per run, acc $0.902 \pm 0.008$ ($+2.6\%$)
\end{itemize}

\textbf{However}: on CIFAR-100, $\Delta$ shrinks to $+0.0006$ (equivalent). The finding is regime-specific: when loss curve is too smooth to plateau, structural signals provide an alternative.

\subsection{Adaptive encoder (Sprint 3)}
Rendering Lorenz trajectory as 64$\times$64 RGB frames, encoding with various models, then analyzing trajectory:

\begin{table}[h]
\centering\small
\begin{tabular}{lrr}
\toprule
encoder & params & corr\_dim error $|$$\Delta_{GT}$$|$ \\
\midrule
raw Lorenz (3D ground truth) & 0 & 0 \\
\textbf{vg\_tiny (random-init)} & \textbf{498} & \textbf{0.247} \\
vg\_small & 6\,362 & 0.740 \\
vg\_medium & 24\,234 & 0.454 \\
ResNet50 (ImageNet pretrained) & 23\,508\,032 & 1.420 \\
\bottomrule
\end{tabular}
\end{table}

\textbf{Random-init vg\_tiny (498 params) is $6\times$ more accurate at recovering corr\_dim than ResNet50 (23.5M params), 10$\times$ faster.} ImageNet-pretrained encoders impose semantic bias that distorts dynamical analysis.

\subsection{Strange loop (Sprint 5C, 5C.1)}
A \texttt{SelfReferentialVG} takes $(x, \text{prev\_state})$ as input where prev\_state = render($\textit{self}$) $\to$ VG-CNN encode. Loss = task + $\alpha \cdot \text{self\_surprise}$ where self\_surprise = MSE(predicted, actual) of current self-state.

With $\alpha=0$: task converges (acc $0.877$), self\_drift $\to 0$ (state stabilizes), but self\_surprise \emph{explodes} from $0.04$ to $2.78$. \emph{System achieves asymmetric steady state: stable self, false self-model, solved task.}

Sweep results (MNIST, 3 seeds):

\begin{table}[h]
\centering\small
\begin{tabular}{rrr}
\toprule
$\alpha$ & MNIST acc & self\_surprise final \\
\midrule
0.0 & $0.877 \pm 0.005$ & $2.540$ \\
0.1 & $0.878 \pm 0.005$ & $0.012$ \\
\textbf{1.0} & $\mathbf{0.881 \pm 0.003}$ & $0.008$ \\
2.0 & $0.880 \pm 0.001$ & $0.007$ \\
5.0 & $0.867 \pm 0.005$ & $0.004$ \\
10.0 & $0.857 \pm 0.011$ & $0.003$ \\
\bottomrule
\end{tabular}
\end{table}

Replicated on CIFAR-100: $\Delta_{\alpha{=}1 \text{ vs } \alpha{=}0} = +1.7\%$ ($4\times$ larger than MNIST $+0.4\%$).

\textbf{Self-coherence is cheap when explicit, and its benefit scales with task difficulty.}

\subsection{Transfer (Sprint 7)}
Pretrain MNIST with $\alpha$, freeze trunk, linear probe FashionMNIST:

\begin{table}[h]
\centering\small
\begin{tabular}{rr}
\toprule
$\alpha$ & Fashion transfer \\
\midrule
0.0 & $0.626 \pm 0.008$ \\
1.0 & $0.627 \pm 0.001$ \\
\textbf{2.0} & $\mathbf{0.635 \pm 0.004}$ \\
\bottomrule
\end{tabular}
\end{table}

Pattern non-monotonic: on MNIST direct $\alpha=1$ best, but for transfer $\alpha=2$ wins. \textbf{Stronger self-coherence pressure breaks task-specificity and produces more transferable features.}

\subsection{Hierarchical observation (Sprint 8C)}
Stacking observers (L1 observes L0, L2 observes L1) does not produce meta-information:

\begin{itemize}
    \item L0 SSVG: $0.077 \to 0.017$ (converges, $-78\%$)
    \item L1 SSVG: $0.057 \to 0.929$ (\textbf{diverges}, $+1531\%$)
    \item L0$\leftrightarrow$L1 Pearson correlation of encoded trajectories: $0.95$--$0.98$
\end{itemize}

Higher-order observers \emph{relay} information without predicting it. Recursive observation is structurally redundant but predictively chaotic. \emph{One observation is enough.}

\subsection{Multi-step horizon (Sprint 8B, null finding)}
Predicting $z[t{+}k]$ for $k \in \{1, 2, 3, 5, 10\}$ produces identical results to single-step. Reason: state stabilizes (drift $\to 0$) under self-coherence loss, so all horizons see indistinguishable inputs.

\section{Discussion}

\subsection{The pattern: less is more}
Three independent findings point in the same direction:
\begin{itemize}
    \item Sprint 3: vg\_tiny (498p) beats ResNet50 (23.5M) on dynamical analysis
    \item Sprint 2: structural signal ($\varphi_{id}$) beats aggregate signal (loss)
    \item Sprint 5C.1: tiny self-coherence pressure ($\alpha{=}0.1$) suffices to eliminate self-model drift
\end{itemize}

Wrong inductive bias dominates capacity. Pretrained ResNet50 sees ImageNet categories in noise; loss aggregates everything including bias; large self-prediction weight ($\alpha{>}5$) becomes anti-task. The right amount of structure beats the maximal amount.

\subsection{Architecture as dynamical system}
Sprint 4 shows: growing networks are not arbitrary fluctuating systems. They have universal spectral signature ($1/f^2$), low-rank dynamics, deterministic-laminar evolution, multi-attractor basins. This is \emph{measurable physics of adaptation}.

\subsection{Self-prediction as IB regularization}
Sprints 5C.1, 6, 7 converge on a clean interpretation: explicit self-prediction loss = information bottleneck pressure. Sweet spot depends on use:
\begin{itemize}
    \item Direct task (any difficulty): $\alpha \approx 1$
    \item Transfer downstream: $\alpha \approx 2$
\end{itemize}

The mechanism: state\_predictor must compress trunk activations into low-D state. This forces structure independent of task signal --- explaining why higher $\alpha$ improves transfer at slight task cost.

\subsection{Gödel-Tarski operational}
Sprints 5C + 8C show concretely that systems cannot self-model consistently:
\begin{itemize}
    \item Internal self-reference (5C): self\_surprise diverges by default
    \item External recursive observation (8C): higher-order observers fail to predict
\end{itemize}

\emph{The structure that is measurable from outside (Sprint 4: spectral, $\varphi_{id}$, basins) is not accessible from inside.} This is the Gödel limit instantiated in code. Self-coherence (Sprint 5C.1) is achievable only with explicit external pressure --- it does not emerge spontaneously.

\subsection{Why biology rarely evolves it}
If self-coherence is computationally cheap given explicit gradient, why is it rare biologically? Probably because environment does not provide an explicit "did you predict your own state?" gradient. Without that signal, drift dominates by default (Sprint 5C). Computational consciousness is engineering-cheap but biologically expensive because the gradient must be invented.

\section{Limitations}

\begin{itemize}
    \item $n=3$ seeds throughout: directional findings, magnitudes uncertain
    \item MNIST and CIFAR-100 flattened (no convolutional structure exploited)
    \item $\varphi_{id}$ implementation is an approximation, not PyPhi
    \item No comparison against state-of-the-art continual learning baselines for VG growth
    \item Transfer tested on FashionMNIST (close distribution) ; harder transfer (MNIST $\to$ CIFAR) untested
\end{itemize}

\section{Conclusion}

We documented seven concrete findings from coupling adaptive neural networks with trajectory analysis tools. The robust findings:

\begin{enumerate}
    \item Growing architectures follow universal dynamical laws ($1/f^2$ spectral)
    \item Random-init tiny CNNs beat large pretrained encoders on dynamical analysis (semantic bias is a feature, not a bug, only for semantic tasks)
    \item Self-prediction loss acts as IB regularization, optimal weight depends on use ($\alpha=1$ task, $\alpha=2$ transfer)
    \item Self-coherence does not emerge spontaneously but is cheap when explicit
    \item Recursive observation produces redundancy, not meta-information
\end{enumerate}

The framework suggests a path forward: instead of scaling pretrained models, build small adaptive observers measuring their own dynamics, gated by self-prediction regularization. This is engineering, not magic. The interesting question --- can such a system bootstrap its own self-coherence gradient --- remains open.

\bibliographystyle{plain}
\begin{thebibliography}{9}
\bibitem{aglas2026verifygrow} Z.~Aglas. VerifyGrowNet: Disciplined Growth for Adaptive Neural Architectures. Workshop draft, 2026.
\bibitem{aglas2026dynamoscope} Z.~Aglas. Video as Latent State-Space Trajectory. arXiv draft, 2026.
\bibitem{friston} K.~Friston. The free-energy principle. \emph{Nature Reviews Neuroscience}, 2010.
\bibitem{hoel} E.~Hoel. The emergence of consequence: a quantitative framework for causal emergence. \emph{Entropy}, 2017.
\bibitem{tishby} N.~Tishby, F.~Pereira, W.~Bialek. The information bottleneck method. 1999.
\bibitem{hofstadter} D.~Hofstadter. \emph{I Am a Strange Loop}. Basic Books, 2007.
\bibitem{wolfram} S.~Wolfram. \emph{A New Kind of Science}. Wolfram Media, 2002.
\end{thebibliography}

\end{document}
